% =========================================================
% Proof Stub: Finite Modification Theorem
% Source: volume-iii/analysis/sequences/notes/sequences/notes-k-tails.tex
% Mode: standard
% =========================================================

\clearpage
\phantomsection
\hypertarget{proof-RA-ABB-T-finite-modification}{}

\subsubsection[Finite Modification Theorem]{Proof --- Finite Modification Theorem}

\bigskip

\noindent
\textbf{Source.}
See \texttt{notes-k-tails.tex}.

\vspace{0.75em}

\noindent
\textbf{Goal.}
If $(x_n)$ and $(y_n)$ differ in only finitely many terms and $(x_n) \to L$, then $(y_n) \to L$.

\vspace{0.75em}

\noindent
\textbf{Logical form.}
\[
  (x_n \ne y_n \text{ for only finitely many } n) \wedge (x_n \to L) \Rightarrow y_n \to L.
\]

\vspace{0.75em}

\noindent
\textbf{Proof strategy.}
Let $N_0$ be past the last index where they differ. For $n \ge N_0$, $x_n = y_n$, so the $\varepsilon$-condition on $(x_n)$ applies to $(y_n)$ directly.

\vspace{0.75em}

\noindent
\textbf{Proof.}
\begin{proof}
\Claim If $(x_n)$ and $(y_n)$ differ in only finitely many terms and $(x_n) \to L$, then $(y_n) \to L$.

\Given 

\Goal 

\Strategy Let $N_0$ be past the last index where they differ. For $n \ge N_0$, $x_n = y_n$, so the $\varepsilon$-condition on $(x_n)$ applies to $(y_n)$ directly.

\medskip

% --- let N_0 = 1 + max{n : x_n ≠ y_n} (finite set, so max exists) ---
% --- for n ≥ N_0: x_n = y_n ---
% --- given ε > 0, let N = max(N_ε, N_0) where N_ε comes from convergence of (x_n) ---
% --- for n ≥ N: |y_n - L| = |x_n - L| < ε ---

\end{proof}

\begin{remark*}[Proof shape]
Past the last differing index, the two sequences agree, so the $\varepsilon$-$N$ argument for $(x_n)$ applies verbatim.
\end{remark*}

\begin{remark*}[Dependencies]
\begin{itemize}
  \item $\varepsilon$-$N$ definition of convergence.
  \item Finite sets have a maximum (well-ordering of $\mathbb{N}$ or induction).
\end{itemize}
\end{remark*}
